\section{Prompt Engineering: ingeniería de prompts}

\subsection{Role Prompting (prompt de rol)}

Role Prompting es la técnica de prompt más simple pero sumamente eficaz. El docente mostró un ejemplo excelente:

\begin{quote}
\textbf{Prompt 1}: What is $100 \times 100 \div 400 \times 56$? \\
\textbf{Respuesta del modelo}: 280 \quad {\color{red}(incorrecta)}
\end{quote}

\begin{quote}
\textbf{Prompt 2}: You are an MIT mathematician. What is $100 \times 100 \div 400 \times 56$? \\
\textbf{Respuesta del modelo}: 1400 \quad {\color{green!60!black}(correcta)}
\end{quote}

Con solo agregar las palabras «You are an MIT mathematician», el modelo pasó de dar una respuesta incorrecta a dar una respuesta correcta.

La explicación intuitiva del docente es la siguiente:

\begin{importantbox}{Por qué funciona Role Prompting — la condicionalización del espacio de probabilidad}
Los datos de entrenamiento del modelo provienen de internet. En internet, una gran cantidad de usuarios dieron respuestas incorrectas al contestar preguntas de matemáticas. Pero si te limitas a las personas que respondieron preguntas comenzando con «soy un matemático del MIT», la probabilidad de que den la respuesta correcta aumenta considerablemente. En esencia, Role Prompting condiciona (conditioning) la distribución de probabilidad de la salida del modelo, concentrando la masa de probabilidad en las regiones donde se encuentran las respuestas de alta calidad dentro de los datos de entrenamiento.
\end{importantbox}

El docente añadió con humor que, debido a la ambigüedad de los word embeddings, «MIT mathematician» probablemente también activa la región del espacio de embeddings donde se encuentra «Harvard mathematician»; por eso esta técnica también funciona para los matemáticos de Harvard.

\subsection{Chain-of-Thought Prompting (prompt de cadena de pensamiento)}

Chain-of-Thought (CoT) Prompting es una técnica que guía al modelo para que «muestre los pasos para resolver el problema».

\textbf{Prompting estándar}: se da un problema de matemáticas con su respuesta como ejemplo one-shot, y después se plantea un problema nuevo. El modelo genera directamente una respuesta (incorrecta).

\textbf{CoT Prompting}: de la misma manera se da un problema de matemáticas, pero el ejemplo incluye los pasos completos de razonamiento, y después se plantea un problema nuevo. El modelo imita este patrón de «mostrar los pasos», deduce paso a paso antes de dar la respuesta final y, con ello, es más probable que llegue a la respuesta correcta.

\begin{importantbox}{La magia de «Let's think step by step»}
Investigaciones posteriores mostraron un hallazgo sorprendente: ni siquiera hace falta dar un ejemplo completo de razonamiento; basta con agregar al final del prompt la frase «Let's think step by step» para activar el modo de razonamiento paso a paso del modelo y aumentar considerablemente la exactitud en tareas de matemáticas y lógica.
\end{importantbox}

La explicación intuitiva del docente sobre la eficacia de CoT se basa en la perspectiva del \textbf{aprendizaje impulsado por errores} (error-driven learning):

\begin{knowledgebox}{Explicación intuitiva de la eficacia de CoT — el área de superficie de error}
Con el prompting estándar, la salida de un problema de matemáticas podría ser «The answer is 27». En toda esa salida, el único lugar donde el modelo puede equivocarse es el número final «27»; el área de superficie de error (error surface area) es pequeña, y la señal de aprendizaje que el modelo obtiene por retropropagación durante el entrenamiento también es limitada.

Cuando se usa CoT, el modelo debe generar los pasos intermedios completos, como «There are 16 players... half of them quit, so 8 remain... a quarter of those lost, so 2 lost... 8 - 2 = 6». Cada paso intermedio es un posible punto de error. Durante el entrenamiento, esta secuencia de salida más larga ofrece una mayor área de superficie de error, de modo que el modelo tiene más oportunidades de equivocarse y, a partir de eso, actualizar sus pesos, con lo que aprende a razonar con mayor precisión.

Nota: este proceso de aprendizaje ocurre durante el \textbf{entrenamiento}; en el momento de la inferencia, los pesos del modelo no se actualizan. CoT funciona en la inferencia porque el modelo ya adquirió esa capacidad durante el entrenamiento, a partir de datos de este tipo de razonamiento paso a paso.
\end{knowledgebox}

\subsection{Resumen de la sección}

Las dos técnicas centrales de la ingeniería de prompts —Role Prompting y Chain-of-Thought Prompting— consisten, en esencia, en manipular la distribución de probabilidad del modelo. Role Prompting condiciona la distribución de salida mediante la definición de una identidad; CoT amplía el espacio de «pensamiento» del modelo al exigir un razonamiento paso a paso. Ambas técnicas pueden combinarse para obtener el mejor resultado.

